% from ega4-2-fr.tex
\oldpage[IV-2]{87}

\begin{corollary}[5.1.3]
\label{IV.2.5.1.3-fr}
Pour toute partie fermée $Y$ d'un préschéma $X$, on a
\begin{equation}
\operatorname{codim}(Y,X)=\inf_{y\in Y}\dim(\mathcal O_{X,y}).
\tag{5.1.3.1}\label{IV.2.5.1.3.1-fr}
\end{equation}
En outre, si $X$ est localement noethérien, on a, pour tout $x\in Y$
\begin{equation}
\operatorname{codim}_x(Y,X)=
\inf_{\substack{y\in Y,\,x\in\overline{\{y\}}}}
\dim(\mathcal O_{X,y}).
\tag{5.1.3.2}\label{IV.2.5.1.3.2-fr}
\end{equation}
\end{corollary}

\begin{proof}
La relation (5.1.3.2) résulte en effet de (5.1.3.1) et de $(0,14.2.6)$.

Ce corollaire permet de \emph{définir}, pour une partie \emph{quelconque}
$Y$ d'un préschéma $X$, la \emph{codimension} $\operatorname{codim}(Y,X)$ de
$Y$ dans $X$ comme égale au second membre de (5.1.3.1).
\end{proof}

\begin{proposition}[5.1.4]
\label{IV.2.5.1.4-fr}
Pour tout préschéma $X$, on a
\begin{equation}
\dim(X)=\sup_{x\in X}(\dim(\mathcal O_x)).
\tag{5.1.4.1}\label{IV.2.5.1.4.1-fr}
\end{equation}
Si toute partie fermée irréductible de $X$ contient un point fermé, on a aussi
\begin{equation}
\dim(X)=\sup_{x\in F}(\dim(\mathcal O_x))
\tag{5.1.4.2}\label{IV.2.5.1.4.2-fr}
\end{equation}
où $F$ est l'ensemble des points fermés de $X$.
\end{proposition}

\begin{proof}
Il suffit de remarquer (en vertu de (I, 2.4.2)) que les chaînes de parties
fermées irréductibles de $X$ correspondent biunivoquement aux chaînes de
parties fermées irréductibles d'un schéma local de $X$; lorsque toute partie
fermée irréductible de $X$ contient un point fermé, on peut évidemment se
borner aux schémas locaux aux points fermés de $X$.

On notera que toute partie fermée irréductible de $X$ contient un point fermé
lorsque $X$ est \emph{quasi-compact} $(0_{\mathrm I},2.1.3)$; nous verrons un
peu plus loin (5.1.11) qu'il en est de même lorsque $X$ est \emph{localement
noethérien}.
\end{proof}

\begin{corollary}[5.1.5]
\label{IV.2.5.1.5-fr}
Pour qu'un préschéma $X$ soit caténaire, il faut et il suffit que, pour tout
$x\in X$, l'anneau local $\mathcal O_x$ soit caténaire. Si de plus toute
partie fermée irréductible de $X$ contient un point fermé, il suffit, pour que
$X$ soit caténaire, que $\mathcal O_x$ soit caténaire pour tout point fermé
$x$ de $X$.
\end{corollary}

\begin{proof}
Le raisonnement est le même que dans (5.1.4), puisqu'il s'agit de comparer des
chaînes de parties fermées irréductibles ayant mêmes extrémités.
\end{proof}

\begin{env}[5.1.6]
\label{IV.2.5.1.6-fr}
Nous allons maintenant examiner les propriétés plus spéciales liées à des
hypothèses noethériennes.

Rappelons $(0,16.2.3)$ qu'un anneau \emph{local noethérien} $A\ne0$ est de
dimension \emph{finie}, égale au nombre minimum de générateurs d'un idéal de
définition de $A$; pour tout idéal premier $\mathfrak p$ de $A$, la hauteur de
$\mathfrak p$, égale à $\dim(A_{\mathfrak p})$, est donc aussi finie. Ces
propriétés, jointes à (5.1.2), (5.1.3) et (5.1.4), montrent que:
\end{env}

\begin{proposition}[5.1.7]
\label{IV.2.5.1.7-fr}
Pour toute partie fermée non vide $Y$ d'un préschéma localement noethérien
$X$, $\operatorname{codim}(Y,X)$ est finie. Si $X$ est noethérien et affine
et $Y$ une partie fermée irréductible de $X$, $\operatorname{codim}(Y,X)$ est
égal au nombre minimum des sections $s_i$ de $\mathcal O_X$ au-dessus de $X$
telles que $Y$ soit une composante irréductible de l'ensemble des $x\in X$
tels que $s_i(x)=0$ pour tout $i$.
\end{proposition}

\begin{corollary}[5.1.8]
\label{IV.2.5.1.8-fr}
Soient $X$ un préschéma localement noethérien, $\mathcal L$ un
$\mathcal O_X$-Module inversible, $f$ une section de $\mathcal L$ au-dessus
de $X$. Alors toute composante irréductible de l'ensemble $Z$

